% ================================================================
% Insertable course-note section: backward equations and MFPTs
% Compatible with notes.tex (amsmath, amssymb already loaded).
% ================================================================
\section{Backward master equations and mean first-passage times}
\label{sec:backward-mfpt}

In previous sections we used the \emph{master equation} to calculate the
probability $P(n,t)$ of finding a population with $n$ individuals at time $t$.
For example, we calculated stationary probability distributions for
immigration--death processes. However, a different question is often more
relevant to ecology and resource management: \emph{how long will it take a
population to reach a certain abundance for the first time?} For example,
we may want to predict the time to extinction of a small population, or the
time until an exploited fish population falls below a critical threshold.
These are \emph{first-passage} questions.

\subsection{A simple example: time to extinction in a pure-death process}
\label{subsec:pure-death-mfpt}
Consider a population that cannot reproduce or receive immigrants. Each
individual dies independently at a constant per-capita rate $\mu>0$.
If the population contains $n$ individuals, the only possible event is
\begin{equation}
 n\longrightarrow n-1,\qquad r_n=\mu n,\qquad g_n=0.
 \label{eq:mfpt-pure-death-rates}
\end{equation}
The population eventually reaches $n=0$, an absorbing state: once extinction
occurs, no further event can occur. Suppose we start with exactly $n$
individuals. Let $\tau_0$ denote the \emph{first time} the population
reaches zero:
\begin{equation}
 \tau_0=\inf\{t\geq 0:N(t)=0\}.
 \label{eq:mfpt-hitting-time}
\end{equation}
The first-passage time $\tau_0$ is a random variable: different stochastic
realizations become extinct at different times. Our aim is to calculate
\begin{equation}
 T_n=\mathbb{E}[\tau_0\mid N(0)=n],\qquad T_0=0,
 \label{eq:mfpt-definition}
\end{equation}
which is called the \emph{mean first-passage time} (MFPT) to extinction.
Notice that $T_n$ is a function of the \emph{initial state} $n$, not of
elapsed time $t$.

We can calculate $T_n$ without solving the full time-dependent master
equation. If the population currently has $n>0$ individuals, the waiting
time $W$ until the next death is exponentially distributed with rate $\mu n$.
Therefore,
\begin{equation}
 \mathbb{E}[W]=\frac{1}{\mu n}.
\end{equation}
After that event, the system contains $n-1$ individuals. Because the process
is Markovian, the expected \emph{additional} time to extinction is $T_{n-1}$.
By adding the two contributions we obtain
\begin{equation}
 T_n=\frac{1}{\mu n}+T_{n-1},\qquad n\geq1.
 \label{eq:mfpt-pure-death-recursion}
\end{equation}
Applying this relation repeatedly gives a simple exact result:
\begin{equation}
 \boxed{T_n=\frac{1}{\mu}\sum_{k=1}^{n}\frac{1}{k}
       =\frac{H_n}{\mu}},
 \label{eq:mfpt-pure-death-solution}
\end{equation}
where $H_n$ is the $n$th harmonic number. For large $n$,
$H_n\simeq\ln n+\gamma$, where $\gamma$ is Euler's constant. Thus,
$T_n$ grows approximately logarithmically, not linearly, with the initial
population size. Even this elementary example shows how transition rates
allow us to predict a biologically meaningful time scale.

\subsection{General derivation for a one-dimensional birth--death process}
\label{subsec:backward-derivation}
Let a continuous-time Markov process $N(t)$ take values
$n=0,1,2,\ldots$, with the two possible transitions
\begin{equation}
 n\longrightarrow n+1\quad\text{at rate }g_n,
 \qquad
 n\longrightarrow n-1\quad\text{at rate }r_n.
 \label{eq:mfpt-transition-rates}
\end{equation}
The \emph{forward} master equation describes the evolution of the
probability of occupying the destination state $n$:
\begin{equation}
 \frac{dP(n,t)}{dt}=
 g_{n-1}P(n-1,t)+r_{n+1}P(n+1,t)
 -(g_n+r_n)P(n,t),
 \label{eq:mfpt-forward}
\end{equation}
with appropriate modifications at the boundaries.

Now consider a target set $A$ (for example, $A=\{0\}$ for extinction, or
$A=\{0,1,\ldots,n_{\mathrm c}\}$ for collapse below a threshold). Define
$T_n$ as the expected time to first enter $A$, given that the process starts
in state $n$. Clearly,
\begin{equation}
 T_n=0\quad\text{for }n\in A.
 \label{eq:mfpt-target-boundary}
\end{equation}
For $n\notin A$, let $R_n=g_n+r_n$ be the total rate of events. Assuming
$R_n>0$, the waiting time to the next event has average $1/R_n$, while the
next event is a birth with probability $g_n/R_n$ and a death with
probability $r_n/R_n$. Conditioning on that first event yields
\begin{equation}
 T_n=\frac{1}{R_n}
     +\frac{g_n}{R_n}T_{n+1}
     +\frac{r_n}{R_n}T_{n-1}.
 \label{eq:mfpt-first-step}
\end{equation}
Multiplying by $R_n$ and rearranging,
\begin{equation}
 \boxed{-1=g_n\bigl(T_{n+1}-T_n\bigr)
             +r_n\bigl(T_{n-1}-T_n\bigr),
             \qquad n\notin A.}
 \label{eq:mfpt-backward}
\end{equation}
This is the \textbf{backward equation for the mean first-passage time}.
It is called \emph{backward} because its unknown $T_n$ is indexed by the
\emph{initial} state from which we start the process. By contrast, the
forward master equation follows the probabilities of the possible states
\emph{at a later time}. Both are descriptions of the same stochastic process.

More generally, if $q_{nm}$ denotes the transition rate from state $n$ to a
different state $m$, the backward equation reads
\begin{equation}
 -1=\sum_{m\neq n}q_{nm}(T_m-T_n),\qquad n\notin A,
 \label{eq:mfpt-general-generator}
\end{equation}
subject to $T_n=0$ on $A$. In matrix notation, if $Q$ is the generator
of the continuous-time Markov chain (with rows corresponding to initial
states), then the unknown MFPTs outside the target solve
\begin{equation}
 Q_{CC}\,\mathbf T_C=-\mathbf 1,
 \label{eq:mfpt-matrix}
\end{equation}
where $C$ is the set of non-target states. This gives a linear system of
equations rather than a system of differential equations. For a finite
state space, with an almost-surely reachable target and finite mean hitting
times, the relevant linear system has a unique finite solution. For an
infinite state space, additional conditions may be required to select a
finite, physically meaningful solution; a finite MFPT should not be assumed
automatically.

\subsection{A second worked example: immigration, births, and deaths}
\label{subsec:mfpt-threshold-example}
Imagine a small population at an initially empty site. Immigrants arrive
at rate $\lambda>0$; resident individuals reproduce at per-capita rate
$\beta\geq0$ and die at per-capita rate $\mu\geq0$. We ask: \emph{how long
will it take until the population first reaches two individuals?} This
question has a very simple answer even though both upward and downward
transitions are possible. Only the states $n=0,1$ matter before reaching
our target $A=\{2,3,\ldots\}$.

The relevant rates are
\begin{equation}
 g_0=\lambda,\quad r_0=0,\qquad
 g_1=\lambda+\beta,\quad r_1=\mu.
\end{equation}
Our boundary condition is $T_2=0$. The two backward equations are
\begin{align}
 -1&=\lambda(T_1-T_0),\label{eq:mfpt-two-state-0}\\
 -1&=(\lambda+\beta)(T_2-T_1)+\mu(T_0-T_1).
 \label{eq:mfpt-two-state-1}
\end{align}
The first equation gives $T_0=T_1+1/\lambda$. Substituting this and
$T_2=0$ into the second equation yields
\begin{equation}
 \boxed{T_1=\frac{1+\mu/\lambda}{\lambda+\beta},\qquad
 T_0=\frac{1}{\lambda}+\frac{1+\mu/\lambda}{\lambda+\beta}.}
 \label{eq:mfpt-two-state-solution}
\end{equation}
For example, when $\beta=0$, a death from state $1$ forces us to wait for
an additional immigrant before trying again to reach state $2$. This
explains why increasing $\mu$ increases the expected hitting time. Notice
that reaching two individuals is a \emph{first-passage} event; state 2 does
not have to be absorbing in the original population model. We only impose
$T_2=0$ because we stop the clock upon reaching the target.

\subsection{Connection with resource collapse and recovery}
\label{subsec:mfpt-harvesting}
Consider now a resource population (for instance, fish) that reproduces,
experiences density-dependent mortality, and is harvested by consumers.
An illustrative stochastic model may use the transition rates
\begin{equation}
 g_n=\lambda+bn,
 \qquad
 r_n=dn+\frac{b}{K}n(n-1)+\frac{hn}{a+n},
 \label{eq:mfpt-harvest-rates}
\end{equation}
where $\lambda$ is immigration, $b$ and $d$ are intrinsic demographic
parameters, $K$ sets the strength of density dependence, and the final
term is a saturating Holling-type-II harvesting rate. The term involving
$n(n-1)$ ensures that density-dependent competitive deaths occur only
when at least two individuals are present. This is an \emph{illustrative}
choice of microscopic rates, not a claim about the exact transition rates
used in any particular published study.

If $n_{\mathrm c}$ is a critical low-abundance threshold, a collapse time
can be defined as the first time a population starting above that threshold
reaches $n\leq n_{\mathrm c}$. The corresponding MFPT satisfies
Eq.~(\ref{eq:mfpt-backward}) for $n>n_{\mathrm c}$, with
$T_n=0$ for $n\leq n_{\mathrm c}$. Similarly, a recovery time from low
abundance can be defined by choosing an upper target threshold.
For numerical calculations we may truncate the state space at some upper
value $n_{\max}$, but then we must specify what happens at that artificial
boundary (for instance, a reflecting boundary with $g_{n_{\max}}=0$).
Predictions should be checked for sensitivity to this truncation.

These questions are particularly relevant to the analysis of nonlinear
harvesting and stochastic collapse/recovery in the contribution by
Cuenda, Llorente, and Capit\'an (2020). In that application, the main
quantity of interest is not simply a stationary distribution: it is the
expected time to cross a biologically significant threshold.

\subsection{Exercises}
\begin{enumerate}
\item Starting from Eq.~(\ref{eq:mfpt-pure-death-recursion}), calculate
$T_1$, $T_2$, and $T_3$. Why is the increase from $T_2$ to $T_3$ smaller
than the increase from $T_1$ to $T_2$?
\item Derive Eq.~(\ref{eq:mfpt-backward}) alternatively by conditioning on
what happens during a short interval $dt$. Include the probability
$1-(g_n+r_n)dt+o(dt)$ that no event occurs.
\item Use a Gillespie simulation to estimate the mean time to extinction
for a pure-death population. Compare the sample average across many
replicates with Eq.~(\ref{eq:mfpt-pure-death-solution}).
\item For the two-individual threshold example, set $\beta=0$ and simplify
Eq.~(\ref{eq:mfpt-two-state-solution}). Discuss the limits $\mu\to0$ and
$\mu\to\infty$.
\item For the harvesting model in Eq.~(\ref{eq:mfpt-harvest-rates}),
choose a collapse threshold, write its backward recurrence explicitly,
and solve the resulting linear system numerically after choosing an
upper boundary. Compare estimates of the MFPT obtained from simulation.
\end{enumerate}

\medskip
\noindent\textbf{Take-home message.} The forward master equation asks
\emph{``Where is the population likely to be at time $t$?''} The backward
MFPT equation asks \emph{``Starting from population size $n$, how long
will it take to reach a specified state for the first time?''} Both
questions can be answered using the same elementary transition rates.
